\chapter{Conclusion}

This research project highlighted the immense complexity of statically verifying memory safety within the Linux kernel. Analyzing these eBPF verifier vulnerabilities demonstrated how even minor logic errors can completely break the security model of the eBPF subsystem.

\section{CVE-2021-3600}
The analysis of CVE-2021-3600 showed how a seemingly harmless register manipulation during the verifier's fixup pass (\texttt{BPF\_MOV32\_REG}) can create a fatal desynchronization between the verifier's mathematical bounds tracking and the actual runtime execution. This desynchronization allowed us to reliably achieve out-of-bounds reads and writes on the kernel heap. However, as demonstrated during Milestone 4, modern kernel mitigations, specifically the \texttt{alu\_limit} pointer arithmetic sanitization, effectively prevent this logical bug from being exploited by unprivileged users, safely restricting the impact.

\section{CVE-2021-3444}
CVE-2021-3444 is the sister bug of CVE-2021-3600, targeting the \textbf{destination} register truncation in 32-bit modulo instead of the source register. By targeting kernel 4.19.19, which predates the \texttt{alu\_limit} pointer arithmetic sanitization, the bounds confusion was fully exploitable from an unprivileged user. Through a six-phase exploit chain: SLUB heap spraying, OOB read for KASLR bypass, \texttt{value\_size} corruption for extended heap R/W, and \texttt{bpf\_map\_ops} hijacking via ret2usr, we achieved a \textbf{full local privilege escalation from uid=1001 to uid=0 (root)}. This is the only CVE in this project where a complete unprivileged LPE was successfully demonstrated.

\section{CVE-2023-52452}
In CVE-2023-52452, I shifted the focus from classic map bounds confusion to stack frame depth tracking. I successfully achieved a reliable information leak from uninitialized kernel stack memory and induced stack corruption by overlapping BPF subprogram frames. While a full privilege escalation was theoretically possible, practical exploitation was heavily constrained by dynamic kernel safeguards and crashes.

\section{CVE-2024-26589}
CVE-2024-26589 introduced a different class of verifier bug: a missing \texttt{case} in the \texttt{switch} statement of \texttt{adjust\_ptr\_min\_max\_vals()}, which failed to restrict variable-offset pointer arithmetic on \texttt{PTR\_TO\_FLOW\_KEYS}. This trivial omission enabled full OOB read and write on the kernel stack via the \texttt{BPF\_PROG\_TYPE\_FLOW\_DISSECTOR} hook. The vulnerability requires \texttt{CAP\_NET\_ADMIN}, making it relevant for container escape and lateral privilege escalation scenarios rather than classical unprivileged LPE.

\section{Comparisons}

\subsection{CVE Comparison}

\begin{table}[H]
\centering
\small
\resizebox{\textwidth}{!}{
\begin{tabular}{@{}lllll@{}}
\toprule
\textbf{Aspect} & \textbf{CVE-2021-3600} & \textbf{CVE-2021-3444} & \textbf{CVE-2023-52452} & \textbf{CVE-2024-26589} \\
\midrule
Bug class & 32-bit truncation & 32-bit truncation & Stack depth miscalc. & Missing switch case \\
Component & \texttt{fixup\_\allowbreak bpf\_\allowbreak calls} & \texttt{fixup\_\allowbreak bpf\_\allowbreak calls} & \texttt{update\_\allowbreak stack\_\allowbreak depth} & \texttt{adjust\_\allowbreak ptr\_\allowbreak min\_\allowbreak max\_\allowbreak vals} \\
Primitive & OOB map R/W & OOB map R/W & OOB stack read & OOB stack R/W \\
OOB Read & Kernel heap & Kernel heap & Stack only & Kernel stack \\
OOB Write & \checkmark (privileged) & \checkmark (unpriv.) & Stack corruption & \checkmark (kernel stack) \\
Info leak & Kernel heap PTR & Kernel heap PTR & Kernel stack PTR & Kernel stack PTR (KASLR) \\
Severity & HIGH (7.8) & HIGH (7.8) & HIGH (7.8) & HIGH (7.8) \\
Kernel tested & 5.10.15 & 4.19.19 & 6.7.1 & 6.7.1 \\
Required privs & Unprivileged & Unprivileged & CAP\_BPF & CAP\_NET\_ADMIN \\
Unpriv LPE & Not achievable & \textbf{Achieved} & Theoretical & N/A (needs capability) \\
Complexity & Medium & High & Very high & Low \\
\bottomrule
\end{tabular}
}
\caption{Comparison between all analyzed eBPF verifier exploits}
\end{table}

\textbf{Key differences:} CVE-2021-3444 stands out as the only vulnerability where a \textbf{full unprivileged LPE} was achieved. While sharing the same root cause class as CVE-2021-3600 (32-bit truncation in \texttt{fixup\_bpf\_calls}), its exploitation on kernel 4.19.19, which lacks the \texttt{alu\_limit} sanitization, demonstrates the critical importance of defense-in-depth: the same bug class is neutralized on 5.10+ but fully exploitable on 4.19. CVE-2024-26589 offers the simplest exploitation path (a 4-line patch fixes it) but requires \texttt{CAP\_NET\_ADMIN}. CVE-2023-52452 has the highest exploitation complexity, targeting the stack frame depth tracking subsystem.

\subsection{Experimental Testing Matrix}

\begin{table}[H]
\centering
\small
\begin{tabularx}{\textwidth}{@{}llXX@{}}
\toprule
\textbf{CVE} & \textbf{Milestone} & \textbf{Expected Result} & \textbf{Observed Result} \\
\midrule
CVE-2021-3600 & M1 & Bounds mismatch visible & Confirmed \\
CVE-2021-3600 & M2 & OOB read primitive & Confirmed (15/15 reads) \\
CVE-2021-3600 & M3 & OOB write primitive & Confirmed (5/5 reliability) \\
CVE-2021-3600 & M4 & Unprivileged full LPE & Blocked by \texttt{alu\_limit} \\
\midrule
CVE-2021-3444 & M1 & Bounds mismatch visible & Confirmed \\
CVE-2021-3444 & M2 & OOB read (KASLR bypass) & Confirmed (\texttt{array\_map\_ops} leaked) \\
CVE-2021-3444 & M3 & OOB write (value\_size corruption) & Confirmed (256$\rightarrow$0x800) \\
CVE-2021-3444 & M4 & Full unprivileged LPE & \textbf{Achieved (uid=1001$\rightarrow$uid=0)} \\
\midrule
CVE-2023-52452 & M1 & Stack data disclosure & Confirmed (kernel pointers leaked) \\
CVE-2023-52452 & M2 & KASLR bypass via leak & Confirmed \\
CVE-2023-52452 & M3 & Type confusion from stack overwrite & Rejected (verifier keeps scalar) \\
CVE-2023-52452 & M4 & Controlled escalation feasibility & Partial: controlled corruption, no full LPE chain \\
\midrule
CVE-2024-26589 & M1 & Verifier bypass on flow\_keys & Confirmed (variable-offset accepted) \\
CVE-2024-26589 & M2 & OOB read beyond flow\_keys & Confirmed (stack data leaked) \\
CVE-2024-26589 & M3 & Controlled OOB with map exfil & Confirmed (KASLR leak) \\
CVE-2024-26589 & M4 & OOB write / full analysis & Confirmed (stack corruption) \\
\bottomrule
\end{tabularx}
\caption{Experimental matrix: expected vs observed outcomes for analyzed milestones}
\end{table}

This matrix provides a compact view of the empirical results and complements the detailed per-CVE exploitation logs.

% \subsection{CVE-2023-52452 vs CVE-2023-2163}

% \begin{table}[H]
% \centering
% \begin{tabular}{@{}lll@{}}
% \toprule
% \textbf{Aspect} & \textbf{CVE-2023-2163} & \textbf{CVE-2023-52452 (ours)} \\
% \midrule
% Bug class & Branch analysis error & Stack depth tracking \\
% Type confusion & Direct (register mismatch) & Indirect (stack overlap) \\
% Arbitrary R/W & Via corrupted stack ptr & Partial (stack corruption) \\
% Full LPE & Achieved & Theoretical \\
% Complexity & High & Very high \\
% \bottomrule
% \end{tabular}
% \caption{Comparison of eBPF verifier stack exploits}
% \end{table}

% CVE-2023-2163's exploit uses a \texttt{CORRUPT\_R6} macro that creates type confusion through carefully crafted conditional branches. The verifier misclassifies a register, allowing packet-controlled data to be used as a stack pointer. This gives direct arbitrary R/W.

% My CVE provides stack corruption rather than direct type confusion. This makes the exploitation path longer: instead of direct R/W, I must corrupt specific kernel stack data and rely on how the kernel uses it. However, the information leak primitive is stronger and more reliable.

\section{Limits and Failed Escalation Summary}

Across all four CVEs, the strongest practical lesson is that proving a primitive and achieving stable unprivileged LPE are different milestones.

\textbf{CVE-2021-3600 limits}
\begin{itemize}
	\item OOB read/write primitives are valid in privileged contexts.
	\item Unprivileged exploitation is blocked by layered mitigations, especially \texttt{alu\_limit} pointer sanitization.
	\item Therefore, this CVE alone does not yield a robust unprivileged root escalation in the tested setup.
\end{itemize}

\textbf{CVE-2021-3444 limits}
\begin{itemize}
	\item Full unprivileged LPE was achieved on kernel 4.19.19, but requires SMEP/SMAP/KPTI to be disabled for the ret2usr technique.
	\item The SLUB heap layout is non-deterministic: identifying the correct file descriptor for the corrupted map required iterating over all candidates.
	\item On kernels with \texttt{alu\_limit} (5.x+), the same vulnerability class is neutralized, demonstrating the effectiveness of defense-in-depth.
\end{itemize}

\textbf{CVE-2023-52452 limits}
\begin{itemize}
	\item Stack leak and controlled corruption are reproducible and technically strong.
	\item Direct type-confusion dereference paths are rejected by verifier checks.
	\item Full escalation would require additional chaining steps (precise stack mapping, controlled fake object placement, and execution redirection).
\end{itemize}

\textbf{CVE-2024-26589 limits}
\begin{itemize}
	\item Full OOB read/write primitives on kernel stack are confirmed and reproducible.
	\item The vulnerability is not reachable from unprivileged users (requires \texttt{CAP\_NET\_ADMIN}).
	\item Full LPE via ROP chain is theoretically achievable but requires precise stack layout knowledge that varies with compiler optimizations and kernel call paths.
	\item Primary real-world impact: container escape when \texttt{CAP\_NET\_ADMIN} is granted.
\end{itemize}

\textbf{Takeaway}
\begin{itemize}
	\item Documenting failed escalation attempts is essential for scientific rigor.
	\item A CVE can be highly relevant even when full LPE is not reached, if the primitive is well-characterized and repeatable.
	\item Different BPF program types expose different attack surfaces; the threat model varies based on required capabilities.
\end{itemize}

\section{Final Thoughts}
The four vulnerabilities analyzed in this report represent distinct but equally critical classes of verifier bugs: post-verification instruction rewriting flaws (CVE-2021-3600 and CVE-2021-3444), dynamic stack depth miscalculations (CVE-2023-52452), and missing pointer type restrictions (CVE-2024-26589). Together, they demonstrate that verifier safety failures span multiple subsystems and BPF program types. The successful full LPE achieved through CVE-2021-3444 on kernel 4.19.19 underscores how the absence of defense-in-depth mechanisms (such as \texttt{alu\_limit}) can elevate a verifier logic bug from a theoretical concern to a practical root compromise. The research highlights that while the eBPF subsystem is a powerful tool for kernel-level extensibility, its security heavily relies on the correctness of the verifier's complex logic. While mitigations like \texttt{unprivileged\_bpf\_disabled=1} drastically reduce the attack surface for regular users, the widespread adoption of eBPF in privileged containers and cloud-native environments means that verifier logic flaws will continue to be a high-value target for kernel exploitation.
